\VerdierPhysicalPageAnchor{0030}%
\VerdierSemanticAnchor{intro-8-local-invariants}
where the $\chi_P(\Phi)$ are \emph{local invariants}, which can be determined
when the sheaf $F$ and the endomorphism $\Phi$ are known in a neighborhood
of $P$. When $X$ is a topological manifold and the sheaf $F$ is constant
and free of rank 1, one recovers the usual Lefschetz formula.

\VerdierSemanticAnchor{intro-8-coherent-case}
In the setting of coherent-sheaf cohomology, there is an analogous formula.
One must take $X$ to be a scheme of finite type and proper over a field $k$,
and $F$ to be a coherent sheaf of finite Tor-dimension, or more generally a
perfect complex. The local invariants that then arise have been studied only
in certain special cases. When the field $k$ is algebraically closed, the
scheme $X$ is smooth over $k$, the sheaf $F$ is locally free, and the fixed
point $P$ is transverse, one obtains
\VerdierSemanticAnchor{intro-8-normal-crossing-formula}
\[
\chi_P(\Phi)
=\frac{\operatorname{Tr}(\Phi_{F,P})}
       {\det\bigl(1-d\Phi_{X,P}\bigr)}
\]

\VerdierSemanticAnchor{intro-8-normal-crossing-data}
Here $\Phi_{F,P}$ denotes the endomorphism induced by $\Phi$ on the reduced
fiber of the sheaf $F$ at $P$ (which is a finite-dimensional $k$-vector
space), and $d\Phi_{X,P}$ denotes the differential of the morphism $\Phi_X$
at $P$.

\VerdierSemanticAnchor{intro-8-residue-extension}
When the fixed point $P$ is no longer assumed to be transverse, i.e. when
the endomorphism $1-d\Phi_{X,P}$ is no longer assumed to be invertible, the
local invariants can still be determined by a residue calculation in the
sense of [4].

\VerdierSemanticAnchor{intro-8-relative-lefschetz}
The Lefschetz theorem in the context of algebraic geometry can moreover be
generalized to relative situations.
\VerdierSemanticAnchor{intro-8-relative-setup}
Let $X\xrightarrow{f}Y$ be a $Y$-scheme proper and smooth over a locally
noetherian scheme $Y$, and let $F$ be a coherent sheaf of finite
Tor-dimension on $X$. The complex of sheaves $\mathrm{R}f_*(F)$ on $Y$ is
then a \emph{perfect complex}. Let
$\Phi\colon(X,F)\longrightarrow(X,F)$ be a $Y$-endomorphism of the scheme
$X$ equipped with the sheaf $F$, and suppose that the scheme $P$ of fixed
points of $\Phi$ is finite over $Y$. The endomorphism $\Phi$ induces an
endomorphism
\VerdierSemanticAnchor{intro-8-relative-rf-phi}
\[
\mathrm{R}f_*(\Phi)\colon
\mathrm{R}f_*(F)\longrightarrow\mathrm{R}f_*(F),
\]

\VerdierSemanticAnchor{intro-8-relative-lefschetz-number}
and since the complex $\mathrm{R}f_*(F)$ is \emph{perfect}, one can define
the Lefschetz number $\chi(\mathrm{R}f_*(\Phi))$, which here is a section of
the structure sheaf of $Y$. With this understood, one has the equality
\VerdierSemanticAnchor{intro-8-relative-formula}
\[
\chi\bigl(\mathrm{R}f_*(\Phi)\bigr)=\chi_P(\Phi),
\]
where $\chi_P(\Phi)$ is a local invariant, which can again be determined by
a residue calculation.

